% =====================================================================
% Figure contract
% 核心结论：CUA 评测必须同时区分评测对象的难度层级与评测方法学的可信度。
% 图原型：mode C / 中等复杂度双轨结构图；上轨含平台 fan-out/fan-in，下轨为方法链。
% 证据层级：hero = 三平台正交汇聚到 Hybrid；支撑 = 对象轴两端演进、方法学链与底部注释。
% 能不能更少？：不能；删任一节点会遗漏用户指定的 §8 章节映射。
%
% Step ① 设计文档（Form A / 中等档）
% 0 参考：沿用 cua-survey-fig1-overview.tex 的 XeLaTeX preamble、低饱和配色与线宽节奏。
% 1 领域：Computer-Use Agents 中文综述；技术术语保持英文。
% 2 格式：TikZ standalone，fontspec + xeCJK，XeLaTeX 编译。
% 3 画布：W_ink = 23.0 cm，含 8 pt border 后仍小于 24 cm；目标印刷宽度 23 cm。
% 4 模块：上轨 7 个阶段（其中平台阶段为 3 张纵排卡）；下轨 3 个方法学节点；底部注释条。
% 5 连线：对象轴实线箭头；平台 fan-out/fan-in 为 trunk + vertical spine + horizontal stubs；
%   Verifiers 到 Hybrid 为单条垂直虚线。
% 6 空间：两条等宽浅底 zone；跨轨标注只占两 zone 之间的空白走廊。
% 7 强调：平台卡和 Hybrid 使用较深描边；其余节点使用低饱和 pastel。
% 8 密度：中等档，无数值或矩阵，不嵌入无依据的 chart。
% 9 ASCII：
%   [静态]→[离线]→┬[Web]────┬→[Hybrid]→[长程]→[个性化·安全·效率]
%                  ├[Mobile]─┤
%                  └[Desktop]┘
%                         ↑ verifier 可靠性
%              [Metrics]→[Verifiers]→[可复现性与污染]
%   [────────────────────── 底部整宽注释条 ──────────────────────]
% 10 可视化决策：输入没有定量数据；正交 fan-in 本身是本图的可记忆结构。
%
% Pre-flight：P1 设计文档已写；P2 同行节点按 positioning 构造；P3 所有端点使用 node anchor；
% P4 N/A；P5 fan-out/fan-in spine 与跨轨走廊已分配；P6 固定 23 cm zone；P7 已读 lessons。
% =====================================================================

\documentclass[border=8pt]{standalone}
\usepackage{fontspec}
\usepackage{xeCJK}
\usepackage{xcolor}
\usepackage{tikz}
\setCJKmainfont{Heiti SC}
\setCJKsansfont{Heiti SC}
\usetikzlibrary{arrows.meta,bending,positioning,calc}

% Low-saturation palette matched to the survey overview figure.
\definecolor{ink}{HTML}{27364A}
\definecolor{muted}{HTML}{748092}
\definecolor{flow}{HTML}{526A80}
\definecolor{zoneBlue}{HTML}{F3F7FA}
\definecolor{zonePurple}{HTML}{F7F4F8}
\definecolor{blueLine}{HTML}{6D879E}
\definecolor{blueFill}{HTML}{E7EEF4}
\definecolor{tealLine}{HTML}{648D8A}
\definecolor{tealFill}{HTML}{E5F0EE}
\definecolor{purpleLine}{HTML}{817594}
\definecolor{purpleFill}{HTML}{EEEAF2}
\definecolor{goldLine}{HTML}{9B875B}
\definecolor{goldFill}{HTML}{F2EEE3}
\definecolor{roseLine}{HTML}{9A777B}
\definecolor{roseFill}{HTML}{F3E9EA}
\definecolor{greenLine}{HTML}{71896E}
\definecolor{greenFill}{HTML}{EAF0E8}
\definecolor{noteFill}{HTML}{F1F2F4}
\definecolor{noteLine}{HTML}{B9C0C8}

\tikzset{
  every node/.style={
    font=\sffamily\fontsize{8.6}{10.4}\selectfont,
    text=ink
  },
  figuretitle/.style={
    anchor=north west,
    font=\sffamily\fontsize{12.6}{15.0}\selectfont\bfseries,
    text=ink
  },
  trackzone/.style={
    rectangle,
    rounded corners=6pt,
    minimum width=23.0cm,
    inner sep=0pt,
    line width=0.75pt
  },
  objectzone/.style={
    trackzone,
    minimum height=5.10cm,
    draw=blueLine!60,
    fill=zoneBlue
  },
  methodzone/.style={
    trackzone,
    minimum height=2.85cm,
    draw=purpleLine!55,
    fill=zonePurple
  },
  zonetitle/.style={
    anchor=north west,
    font=\sffamily\fontsize{9.5}{11.3}\selectfont\bfseries
  },
  objectbox/.style={
    rectangle,
    rounded corners=4pt,
    minimum height=1.15cm,
    align=center,
    inner sep=4.5pt,
    line width=0.75pt,
    font=\sffamily\fontsize{8.5}{10.2}\selectfont
  },
  platformcard/.style={
    rectangle,
    rounded corners=3pt,
    minimum width=3.75cm,
    minimum height=0.80cm,
    align=center,
    inner sep=3.5pt,
    line width=0.70pt,
    fill=blueFill,
    draw=blueLine,
    font=\sffamily\fontsize{8.2}{9.8}\selectfont
  },
  methodbox/.style={
    rectangle,
    rounded corners=4pt,
    minimum height=1.05cm,
    align=center,
    inner sep=5pt,
    line width=0.75pt,
    font=\sffamily\fontsize{8.7}{10.5}\selectfont
  },
  notebar/.style={
    rectangle,
    rounded corners=4pt,
    draw=noteLine,
    fill=noteFill,
    line width=0.6pt,
    minimum width=23.0cm,
    minimum height=1.35cm,
    align=center,
    inner sep=5pt,
    font=\sffamily\fontsize{8.3}{10.2}\selectfont\bfseries
  },
  arrowshort/.style={
    -{Stealth[length=3pt 1.0,width'=0pt 0.5,sep=0pt 0.3]},
    line width=0.85pt,
    shorten >=1pt,
    shorten <=0pt,
    color=flow
  },
  fanstubshort/.style={
    -{Stealth[length=3pt 1.0,width'=0pt 0.5,sep=0pt 0.3]},
    line width=0.85pt,
    shorten >=1pt,
    shorten <=0pt,
    color=flow
  },
  flowline/.style={
    line width=0.85pt,
    color=flow
  },
  crosslink/.style={
    -{Stealth[length=4pt 1.1,width'=0pt 0.55,sep=0pt 0.4]},
    densely dashed,
    line width=0.8pt,
    shorten >=2pt,
    shorten <=1pt,
    color=purpleLine
  },
  crosslabel/.style={
    anchor=east,
    font=\sffamily\fontsize{8.0}{9.6}\selectfont,
    text=purpleLine,
    inner sep=1pt
  }
}

\begin{document}
\sffamily
\begin{tikzpicture}

% Frame and zones.
\node[figuretitle] (title) at (0,0)
  {CUA 评测体系：对象轴 × 方法学轴（§8）};
\node[objectzone,below=0.38cm of title.south west,anchor=north west] (objectZone) {};
\node[methodzone,below=0.72cm of objectZone.south,anchor=north] (methodZone) {};
\node[notebar,below=0.42cm of methodZone.south] (noteBar)
  {\mbox{高 grounding accuracy 推不出 long-horizon completion；}\\[1pt]
   \mbox{横向比较必须同时报告 harness·step budget·backbone·verifier·cost（§1.2、§8.14）}};

% Zone titles.
\node[zonetitle,text=blueLine,
      below right=0.30cm and 0.48cm of objectZone.north west] (objectTitle)
  {评测对象轴};
\node[zonetitle,text=purpleLine,
      below right=0.30cm and 0.48cm of methodZone.north west] (methodTitle)
  {评测方法学轴};

% Upper track: object axis.
\node[objectbox,anchor=west,right=0.65cm of objectZone.west,
      minimum width=2.55cm,text width=2.15cm,fill=blueFill,draw=blueLine] (static)
  {静态 Grounding\\§8.1};
\node[objectbox,right=0.48cm of static,
      minimum width=2.35cm,text width=1.95cm,fill=tealFill,draw=tealLine] (offline)
  {离线轨迹\\§8.2};

\node[platformcard,right=1.05cm of offline] (mobile)
  {Mobile：MobileWorld §8.4};
\node[platformcard,above=0.18cm of mobile] (web)
  {Web：WebArena §8.3};
\node[platformcard,below=0.18cm of mobile] (desktop)
  {Desktop：OSWorld 2.0 §8.5};

\node[objectbox,right=1.05cm of mobile,
      minimum width=2.35cm,text width=1.95cm,fill=purpleFill,draw=purpleLine] (hybrid)
  {Hybrid 接口\\§8.6};
\node[objectbox,right=0.50cm of hybrid,
      minimum width=3.25cm,text width=2.85cm,fill=goldFill,draw=goldLine] (long)
  {长程/专业：\\Odysseys §8.7};
\node[objectbox,right=0.45cm of long,
      minimum width=3.55cm,text width=3.15cm,fill=roseFill,draw=roseLine] (personal)
  {个性化·安全·效率\\§8.8–8.10};

% Main object-axis arrows.
\draw[arrowshort] (static.east) -- (offline.west);
\draw[arrowshort] (hybrid.east) -- (long.west);
\draw[arrowshort] (long.east) -- (personal.west);

% Orthogonal fan-out: offline trajectory to the three platform cards.
\coordinate (fanOutTop) at ([xshift=-0.52cm]web.west);
\coordinate (fanOutMid) at (fanOutTop |- mobile.west);
\coordinate (fanOutBottom) at (fanOutTop |- desktop.west);
\draw[flowline] (offline.east) -- (fanOutMid);
\draw[flowline] (fanOutTop) -- (fanOutBottom);
\draw[fanstubshort] (fanOutTop) -- (web.west);
\draw[fanstubshort] (fanOutMid) -- (mobile.west);
\draw[fanstubshort] (fanOutBottom) -- (desktop.west);

% Orthogonal fan-in: the three platform cards merge before Hybrid.
\coordinate (fanInTop) at ([xshift=0.52cm]web.east);
\coordinate (fanInMid) at (fanInTop |- mobile.east);
\coordinate (fanInBottom) at (fanInTop |- desktop.east);
\draw[flowline] (web.east) -- (fanInTop);
\draw[flowline] (mobile.east) -- (fanInMid);
\draw[flowline] (desktop.east) -- (fanInBottom);
\draw[flowline] (fanInTop) -- (fanInBottom);
\draw[fanstubshort] (fanInMid) -- (hybrid.west);

% Lower track: methodology axis, aligned by construction around Verifiers.
\node[methodbox,minimum width=3.20cm,text width=2.80cm,
      fill=purpleFill,draw=purpleLine] (verifiers)
  at (hybrid |- methodZone.center) {Verifiers §8.12};
\node[methodbox,left=1.20cm of verifiers,
      minimum width=3.20cm,text width=2.80cm,fill=blueFill,draw=blueLine] (metrics)
  {Metrics §8.11};
\node[methodbox,right=1.20cm of verifiers,
      minimum width=5.05cm,text width=4.65cm,fill=greenFill,draw=greenLine] (repro)
  {可复现性与污染 §8.13};

\draw[arrowshort] (metrics.east) -- (verifiers.west);
\draw[arrowshort] (verifiers.east) -- (repro.west);

% Cross-axis reliability link; label occupies the whitespace between zones.
\draw[crosslink] (verifiers.north) -- (hybrid.south);
\coordinate (gapMid) at ($(objectZone.south)!0.5!(methodZone.north)$);
\node[crosslabel] at ([xshift=-0.25cm]verifiers |- gapMid)
  {verifier 可靠性决定对象轴分数的可信度};

\end{tikzpicture}
\end{document}
